%% applications -- subsumption of constrained type schemes, correctness of
%% simplification, and examples: identity schemes, Pottier's missed
%% simplification, and a merge of self-application types that is incorrect.
%% The proof of the subsumption formula: appendix-applications.tex.

\providecommand{\parity}{\operatorname{par}}

\section{Constrained Types and Simplification}\label{sec:ap}

This section states the subsumption of constrained type schemes as an
existential entailment, discusses the correctness of simplifications, and
ends with examples, all of which involve variables that occur both co- and
contravariantly. The proof of the subsumption formula is in \cref{app:ap}.

The signature is $\BFT$ with an arity $n$ and a variance $c$, or $\BF$ with $f$
covariant, and $\mathcal D$ is a tree class. For arrow types, $n=2$, $c(0)=1$,
$c(1)=0$, and we write $a\to b$ for $f(a,b)$. Constraints are inequalities
between finite terms; by \cref{lem:bg-flattening} they are flattened, with the
names of the subterms of the antecedent visible and those of the consequent
hidden. Plain entailments are taken in the same signature and tree class.

\subsection{Subsumption}\label{sec:ap-subsumption}

A \emph{type scheme} $S=\forall\bar a.\,C\Rightarrow(\Gamma\Rightarrow\tau)$
consists of a finite set $\bar a$ of bound type variables, a finite
conjunction $C$ of inequalities between terms, a \emph{context} $\Gamma$, that
is, a finite map from program variables to terms, and a term $\tau$. Its other
type variables are \emph{free}. For a valuation $\eta$ of the free variables
in $\mathcal D$, the \emph{instances} of $S$ under $\eta$ are the pairs
$(\rho\circ\Gamma,\rho(\tau))$, where $\rho$ extends $\eta$ to $\bar a$ with
values in $\mathcal D$, satisfies $C$, and $\rho\circ\Gamma$ maps each $v$ in
$\operatorname{Dom}\Gamma$ to $\rho(\Gamma(v))$. Following
\citet[Section~4.2]{TrifonovSmith96}, an instance $(\Delta,T)$ is below an
instance $(\Delta',T')$ if $T\le_c T'$,
$\operatorname{Dom}\Delta\subseteq\operatorname{Dom}\Delta'$, and
$\Delta'(v)\le_c\Delta(v)$ for every $v\in\operatorname{Dom}\Delta$; contexts
are compared contravariantly. Then $S_1\le S_2$ if, under every valuation of
the free variables, every instance of $S_2$ is above some instance of $S_1$.

\begin{proposition}[Subsumption formula]\label{prop:ap-subsumption}
Let $S_i=\forall\bar a_i.\,C_i\Rightarrow(\Gamma_i\Rightarrow\tau_i)$ for
$i=1,2$, where the bound variables are renamed so that $\bar a_1$ and $\bar a_2$
are disjoint from each other and from the free variables of both schemes. Let
the visible variables be $\bar a_2$ and the free variables, and let $\bar a_1$
be hidden.
\begin{enumerate}[leftmargin=*]
\item If $\operatorname{Dom}\Gamma_1\subseteq\operatorname{Dom}\Gamma_2$, then
  \[
   S_1\le S_2\iff C_2\models_{\mathcal D}\exists\bar a_1\,(C_1\wedge R_{12}),
   \qquad
   R_{12}=(\tau_1\le\tau_2)\wedge
   \bigwedge_{v\in\operatorname{Dom}\Gamma_1}\Gamma_2(v)\le\Gamma_1(v).
  \]
\item If $\operatorname{Dom}\Gamma_1\not\subseteq\operatorname{Dom}\Gamma_2$,
  then $S_1\le S_2$ iff $C_2$ has no solution in $\mathcal D$.
\end{enumerate}
\end{proposition}

Both renamings are needed. For $S_1=b$, with $b$ free, and $S_2=\forall b.\,b$,
with empty contexts, the formula read without renaming $\bar a_2$ would be the
valid entailment $\mathrm{true}\models b\le b$; but at $b=\Top$ the instance
$\Bot$ of $S_2$ lies above no instance of $S_1$, so $S_1\not\le S_2$.

This is \cref{eq:intro-subsumption} with contexts. The formula has no
restriction on the occurrences of variables: in a scheme for a higher-order
program, a type variable typically occurs in the domain and in the codomain of
arrows, and then at both polarities.

\begin{corollary}[Deciding subsumption]\label{cor:ap-decidable}
For $\BFT$ with any arity and variance, in particular for arrow types over
$\Bot$, $\to$ and $\Top$, and for $\BF$ with $f$ covariant, subsumption of type
schemes is decidable in each tree class. Over regular trees, the setting of
\citet{TrifonovSmith96}, it agrees with subsumption over arbitrary trees.
\end{corollary}

\begin{proof}
By \cref{prop:ap-subsumption} and \cref{lem:bg-flattening}, subsumption is an
existential entailment between flat systems, or the test that $C_2$ is
unsatisfiable. The latter is the entailment
$C_2\models_{\mathcal D}f(\Bot,\ldots,\Bot)\le\Bot$, since its consequent is
false in every tree class. After flattening, both entailments are decided by
\cref{thm:as-main}. Regular and arbitrary trees agree for both questions
(\cref{thm:tc-regular}).
\end{proof}

\subsection{Simplification}\label{sec:ap-simplification}

A system $C(\bar x,\bar h)$ may be replaced by $C'(\bar x,\bar k)$, where
$\bar x$ remain observable, if $\exists\bar h\,C\iff\exists\bar k\,C'$. The
forward implication is the existential entailment $C\models\exists\bar k\,C'$
with $\bar h$ visible, and the reverse one exchanges the systems; equivalence
of type schemes is subsumption in both directions
(\cref{prop:ap-subsumption}) and need not preserve the solutions on every type
variable. These are the correctness conditions of constraint simplification
and minimization \citep{Pottier96,Pottier01,TrifonovSmith96,EifrigSmithTrifonov95},
and by \cref{thm:as-main} they are decidable, for arrow types without any
restriction on polarity.

Informally, some simplifications need only plain entailment, such as
deleting a literal that the remaining literals entail. Existential entailment
is needed when a transformation changes the hidden variables: eliminating or
merging internal variables, or replacing a scheme by one with different bound
variables (\cref{ex:ap-selfapp}).

\subsection{Examples}\label{sec:ap-examples}

\begin{example}[Two identity schemes]\label{ex:ap-identity}
Let $S_1=\forall\alpha.\,\alpha\to\alpha$ and
$S_2=\forall\beta\gamma.\,(\beta\le\gamma)\Rightarrow\beta\to\gamma$, with empty
contexts. By \cref{prop:ap-subsumption}, $S_1\le S_2$ iff
$\beta\le\gamma\models\exists\alpha\,(\alpha\to\alpha\le\beta\to\gamma)$. Both
sides of the inequality are arrows, so it says $\beta\le_c\alpha$ and
$\alpha\le_c\gamma$, and $\alpha=\beta$ is a witness. Conversely, $S_2\le S_1$
iff $\models\exists\beta\gamma\,(\beta\le\gamma\wedge
\beta\to\gamma\le\alpha\to\alpha)$, with $\alpha$ visible, that is, iff some
$\beta,\gamma$ satisfy $\alpha\le\beta\le\gamma\le\alpha$, and
$\beta=\gamma=\alpha$ is a witness. So the two schemes are equivalent. The
hidden variable $\alpha$ of the first entailment occurs in the domain and in
the codomain of the same arrow, so it has both polarities, and after
flattening the literal $f(\alpha,\alpha)\le u$ admits no parity assignment
(\cref{rem:ch-parity}), which would need $\pi(\alpha)=\pi(u)\oplus1$ and
$\pi(\alpha)=\pi(u)$.
\end{example}

\citet{TrifonovSmith96} decide subsumption by a sound approximation through
canonical bounds, and simplifiers prove entailments by syntactic reasoning
about the bounds of variables; such tests are incomplete \citep{Pottier01}. The
following example of \citet[Section~8.5]{Pottier98thesis} shows a correct
simplification that reasoning through lower bounds misses.

\begin{example}[A missed redundant constraint]\label{ex:ap-pottier}
For arrow types put
\[
 g=\Bot\to\Top,\qquad k=g\to\Bot,\qquad C(a)=(a\to\Bot\le a).
\]
Then $C\models k\le a$ in each tree class. Indeed $a\ne\Bot$, since
$a\to\Bot\le_c\Bot$ is false. If $a=\Top$, then $k\le_c a$. Otherwise $a$ is an
arrow $a_0\to a_1$, so $a\le_c g$, because $g$ is the greatest arrow
(\cref{ex:lf-arrow}); by contravariance $k=g\to\Bot\le_c a\to\Bot\le_c a$. Thus
adding or deleting the literal $k\le a$ is an exact simplification, and
equivalently $C\models\exists z\,(k\le z\wedge z\le a)$, with the witness
$z=k$. The variable $a$ has the single lower bound $a\to\Bot$
and the upper bound $\Top$ in $C$. Proving $k\le a$ through this lower bound
requires $k\le a\to\Bot$, that is, $a\le g$, and this is not entailed:
$a=\Top$ satisfies $C$. A test that proves an inequality $k\le a$ only through a
lower bound of $a$ therefore fails on this true simplification.

The antecedent mentions $a$ at both polarities, and its flattening admits no
parity assignment. The example also shows why the signed system cannot be read on its
own. The signed literals of $C$ (\cref{tab:bg-signed}) are
$f(a^-,\Bot)\le a^+$ and $a^-\le f(a^+,\Top)$, read covariantly. Put
$F=f(\Bot,\Bot)$. The covariant values $a^+=F$ and $a^-=\Bot$ satisfy them, but
they violate $a^-=(a^+)^\partial$, the condition of \cref{lem:bg-signed}. The
signed form of $k\le a$ requires $\nu_0k\le a^+$
(\cref{lem:bg-normalization}(2)), and $\nu_0k=f(\nu_1g,\Bot)=f(F,\Bot)$ by
\cref{lem:bg-normalization}(1). Since $f(F,\Bot)\not\le F$, a test that drops
the coupling between $a^+$ and $a^-$ reports a countermodel of a true
entailment. The leaf cones of \cref{sec:lf} keep this coupling: the two cones
of $a$ are the cones of one tree, and the exclusions and duality of
\cref{sec:lf-duality} enforce it.
\end{example}

The last example concerns a merge of type variables, which changes the bound
variables of a scheme and so needs existential entailment.

\begin{example}[Merging the variables of self-application]\label{ex:ap-selfapp}
The term $\lambda x.\,x\,x$ has the inferred scheme
\[
 S=\forall\alpha\beta.\,(\alpha\le\alpha\to\beta)\Rightarrow\alpha\to\beta,
\]
in which $\alpha$ is both a function and its own argument. A simplifier might
merge $\beta$ into $\alpha$ and propose
$S'=\forall\gamma.\,(\gamma\le\gamma\to\gamma)\Rightarrow\gamma\to\gamma$.
By \cref{prop:ap-subsumption}, $S\le S'$ iff
\[
 \gamma\le\gamma\to\gamma\ \models\ \exists\alpha\beta\,
 (\alpha\le\alpha\to\beta\ \wedge\ \alpha\to\beta\le\gamma\to\gamma),
\]
which holds with the witness $\alpha=\beta=\gamma$. But $S'\le S$ iff
\[
 \alpha\le\alpha\to\beta\ \models\ \exists\gamma\,
 (\gamma\le\gamma\to\gamma\ \wedge\ \gamma\to\gamma\le\alpha\to\beta),
\]
and the second literal says $\alpha\le_c\gamma\le_c\beta$. The values
$\alpha=\Top\to\Bot$ and $\beta=\Bot$ satisfy the antecedent: the domain $\Top$
of $\alpha$ lies above $\alpha$, and its codomain $\Bot$ below $\beta$. No
$\gamma$ lies between the arrow $\alpha$ and $\Bot$, so the entailment fails, and
the instance $(\Top\to\Bot)\to\Bot$ of $S$, the type of $\lambda x.\,x\,x$ applied
to arguments that accept everything and return $\Bot$, lies above no instance
of $S'$. The merge loses generality. In both entailments every hidden variable
has both occurrence signs in the consequent, and a hidden variable is bounded
by an arrow in which it occurs itself, $\alpha\le\alpha\to\beta$ and
$\gamma\le\gamma\to\gamma$, and \cref{thm:as-main} decides both.
\end{example}
